\subsection{Curvature-Assisted Dynamical Compactification in a Pre-Inflationary Higher-Dimensional Universe --- arXiv:2604.25569}

\textbf{Authors:} Yusuke Yamada

\paragraph{主な主張}
5次元 $S^{1}$ compactification の開いた FRW 宇宙において、負の空間曲率が radion を一時的に freeze させた後に tracker-like evolution を維持し、熱的 KK 補正が薄まって安定化極小が現れるまで overshoot を抑えることで、compactified vacuum への捕獲と後続する4次元 inflation への接続が可能であることを示した \cite{Yamada:2026curvature}。

\paragraph{新規性}
曲率摩擦による scalar の減速自体ではなく、同じ bulk 量子場から late-time の Casimir 安定化 potential と early-time の熱的・KK source を導出し、真に5次元的な膨張相から有効4次元宇宙へ移行できるかを半古典的な連成系として検証した点にある。

\paragraph{位置付け}
静的に安定化済みの余剰次元を初期条件として仮定せず、負曲率が potential を直接変形せずに slowly redshifting な摩擦源として働くことで dynamical compactification を助ける proof of principle であり、string-motivated modulus potential に対する curvature-assisted trapping を検討する際の出発点となる。
